\documentclass[11pt]{article}

\usepackage{amsmath,amssymb,amsthm}
\usepackage{algorithm,algorithmic}
\usepackage{booktabs}
\usepackage{geometry}
\usepackage{hyperref}

\geometry{margin=1in}

\newtheorem{theorem}{Theorem}
\newtheorem{lemma}[theorem]{Lemma}
\newtheorem{proposition}[theorem]{Proposition}
\newtheorem{corollary}[theorem]{Corollary}
\newtheorem{definition}[theorem]{Definition}
\newtheorem{remark}[theorem]{Remark}

\DeclareMathOperator{\diag}{diag}
\DeclareMathOperator{\blockdiag}{block\text{-}diag}
\DeclareMathOperator{\proj}{proj}

\title{BISC: Block Inertia-free Structural Convexification\\[6pt]
\large Minimum-Perturbation Hessian Convexification\\
for OC-Structured Interior Point Methods}

\author{}
\date{}

\begin{document}
\maketitle

\begin{abstract}
We present BISC (Block Inertia-free Structural Convexification), a structure-aware
Hessian convexification method for interior point methods (IPMs) applied to
discrete-time optimal control problems. BISC exploits the block-banded KKT structure
via a backward Riccati recursion, applying the minimum Frobenius-norm positive-definite
projection only to the small control Hessian blocks that are indefinite. The
convexification threshold is tied to the barrier parameter, eliminating free parameters
and trial-and-error factorizations. We prove that BISC produces descent directions,
achieves stagewise minimum-norm perturbation, and recovers the exact Newton step
under second-order sufficient conditions---enabling quadratic convergence when
combined with automatic differentiation for exact Hessians.
\end{abstract}

%=============================================================================
\section{Introduction}
%=============================================================================

Interior point methods for nonlinear optimal control require solving a Newton
system at each iteration. When exact Hessians are available (e.g., from automatic
differentiation), the Hessian of the Lagrangian is generally \emph{indefinite}
due to constraint curvature terms. The Newton step may then fail to be a descent
direction, and the Riccati factorization may break down.

\textbf{Current practice (IPOPT-style):} Attempt a Cholesky factorization of
the reduced Hessian. If the inertia is wrong, add a scalar $\delta$ to the entire
Hessian diagonal: $\tilde{H} \leftarrow \tilde{H} + \delta I$. Retry. Increase
$\delta$. Retry again. This approach is:
\begin{itemize}
  \item \textbf{Wasteful:} modifies all stages, even healthy ones.
  \item \textbf{Slow:} requires multiple trial factorizations (typically 2--5).
  \item \textbf{Over-regularizing:} adds the same $\delta$ everywhere, destroying
        curvature information at stages that were already positive definite.
\end{itemize}

\textbf{BISC:} Exploit the block structure of optimal control. Use a backward
Riccati recursion, and at each stage, apply the \emph{minimum Frobenius-norm
PSD projection} only to the small $n_u \times n_u$ control Hessian block that
actually needs it. One pass, no trial-and-error, minimum perturbation, exact
Newton recovery near the solution.

\textbf{Contributions:}
\begin{enumerate}
  \item A complete IPM formulation with backward Riccati and per-block convexification.
  \item The Cheng--Higham modified Cholesky adapted to the OC block structure.
  \item Proof that BISC gives stagewise minimum Frobenius-norm perturbation.
  \item Proof of exact Newton recovery under SOSC, enabling quadratic convergence
        with AD exact Hessians.
  \item Detailed comparison with Verschueren et al.\ (SQP) and Vanroye et al.\ (dual regularization).
\end{enumerate}


%=============================================================================
\section{Problem Formulation}
%=============================================================================

\subsection{Discrete-Time Optimal Control NLP}

Consider the discrete-time optimal control problem with $N$ stages:
\begin{equation}\label{eq:ocp}
\begin{aligned}
  \min_{x_k, u_k} \quad & \sum_{k=0}^{N-1} \ell_k(x_k, u_k) + \ell_N(x_N) \\
  \text{s.t.} \quad
    & f_k(x_k, u_k) - x_{k+1} = 0, \quad k = 0, \dots, N-1, \\
    & g_k(x_k, u_k) \le 0, \quad k = 0, \dots, N, \\
    & x_0 = \bar{x},
\end{aligned}
\end{equation}
where $x_k \in \mathbb{R}^{n_x}$ is the state, $u_k \in \mathbb{R}^{n_u}$ is
the control (absent at terminal stage $N$), $f_k$ is the dynamics map, $\ell_k$
is the stage cost, and $g_k \in \mathbb{R}^{p_k}$ are inequality constraints.

Define the stage variable $z_k = (x_k, u_k) \in \mathbb{R}^{n_z}$ for
$k = 0, \dots, N-1$ and $z_N = x_N$ at the terminal stage.


\subsection{Interior Point Barrier Reformulation}

Introduce slack variables $s_k \ge 0$ and form the barrier subproblem:
\begin{equation}\label{eq:barrier}
\begin{aligned}
  \min_{z, s} \quad & \sum_k \ell_k(z_k) - \mu \sum_k \mathbf{1}^\top \log(s_k) \\
  \text{s.t.} \quad
    & f_k(x_k, u_k) - x_{k+1} = 0, \quad k = 0, \dots, N-1, \\
    & g_k(z_k) + s_k = 0, \quad k = 0, \dots, N,
\end{aligned}
\end{equation}
where $\mu > 0$ is the barrier parameter, decreased toward zero as the IPM progresses.

The Lagrangian with multipliers $\lambda_k$ (dynamics) and $\nu_k$ (inequalities):
\begin{equation}\label{eq:lagrangian}
  \mathcal{L} = \sum_k \ell_k(z_k) - \mu \sum_k \mathbf{1}^\top \log(s_k)
  + \sum_k \lambda_k^\top \bigl(f_k(x_k, u_k) - x_{k+1}\bigr)
  + \sum_k \nu_k^\top \bigl(g_k(z_k) + s_k\bigr).
\end{equation}

\textbf{Sign convention.} We write the dynamics residual as
$e_k = f_k(x_k, u_k) - x_{k+1} = 0$. The Jacobians are:
\begin{equation}
  \frac{\partial e_k}{\partial x_k} = F_k, \qquad
  \frac{\partial e_k}{\partial u_k} = G_k, \qquad
  \frac{\partial e_k}{\partial x_{k+1}} = -I_{n_x}.
\end{equation}
This convention gives $+G_k^\top \lambda_k$ in the control stationarity equation.


\subsection{Slack Elimination and the Barrier Hessian}

The complementarity condition $S_k \nu_k = \mu \mathbf{1}$ (where $S_k = \diag(s_k)$)
allows elimination of slacks and inequality multipliers. After elimination, the
inequality constraints contribute a barrier Hessian term:
\begin{equation}
  \Sigma_k = \diag\!\left(\frac{\nu_{k,i}}{s_{k,i}}\right)
           = \diag\!\left(\frac{\mu}{s_{k,i}^2}\right) \succeq 0.
\end{equation}


\subsection{The Reduced KKT System}

After slack elimination, the Newton system becomes:
\begin{equation}\label{eq:kkt}
  \begin{bmatrix} \tilde{H} & J^\top \\ J & 0 \end{bmatrix}
  \begin{bmatrix} \Delta z \\ \Delta \lambda \end{bmatrix}
  = -\begin{bmatrix} r_d \\ r_p \end{bmatrix},
\end{equation}
where $\Delta z = (\Delta z_0, \dots, \Delta z_N)$ is the primal step,
$\Delta \lambda = (\Delta \lambda_0, \dots, \Delta \lambda_{N-1})$ is the dual step,
$r_d$ is the dual residual, and $r_p$ is the primal residual.


\subsection{The Augmented Hessian $\tilde{H}$ (Block-Diagonal)}

The augmented Hessian is \emph{block-diagonal}:
\begin{equation}\label{eq:Hblock}
  \tilde{H} = \blockdiag(\tilde{H}_0, \tilde{H}_1, \dots, \tilde{H}_N).
\end{equation}
This holds because $f_k(x_k, u_k)$ depends only on $z_k$, not $z_{k+1}$.
There are no cross-stage second-derivative terms. All cross-stage coupling
comes from the constraint Jacobian $J$.

Each stage block:
\begin{equation}\label{eq:Hk}
  \tilde{H}_k = W_k + J_{g,k}^\top \Sigma_k\, J_{g,k},
\end{equation}
where $W_k$ is the Lagrangian Hessian and the second term is the barrier contribution.

\textbf{Lagrangian Hessian} (from AD, possibly indefinite):
\begin{equation}\label{eq:Wk}
  W_k = \underbrace{\nabla^2_{z_k} \ell_k}_{\text{objective (often PSD)}}
      + \underbrace{\lambda_k^\top \nabla^2_{z_k} f_k}_{\text{dynamics curvature (\textbf{can be negative})}}
      + \underbrace{\nu_k^\top \nabla^2_{z_k} g_k}_{\text{inequality curvature}}.
\end{equation}
The dynamics curvature term $\lambda_k^\top \nabla^2 f_k$ is the source of
indefiniteness. It carries genuine second-order information; dropping it
(Gauss--Newton) limits convergence to linear rate.

\textbf{Barrier contribution} (always PSD, decays with $\mu$):
\begin{equation}
  J_{g,k}^\top \Sigma_k\, J_{g,k} \succeq 0.
\end{equation}

\textbf{Sub-block decomposition:}
\begin{equation}\label{eq:Hk_blocks}
  \tilde{H}_k = \begin{bmatrix} \tilde{Q}_k & \tilde{N}_k \\
                                 \tilde{N}_k^\top & \tilde{R}_k \end{bmatrix},
\end{equation}
where $\tilde{Q}_k \in \mathbb{R}^{n_x \times n_x}$ (state--state),
$\tilde{R}_k \in \mathbb{R}^{n_u \times n_u}$ (control--control),
$\tilde{N}_k \in \mathbb{R}^{n_x \times n_u}$ (cross term).


\subsection{The Constraint Jacobian $J$ (Block-Bidiagonal)}

\begin{equation}\label{eq:J}
  J = \begin{bmatrix}
    F_0 & G_0 & -I &     &     &    \\
        &     & F_1 & G_1 & -I &    \\
        &     &     & \ddots &   & \ddots \\
        &     &     & F_{N-1} & G_{N-1} & -I
  \end{bmatrix}.
\end{equation}
Each row~$k$ encodes the linearized dynamics:
$\Delta x_{k+1} = F_k \Delta x_k + G_k \Delta u_k - r_{p,k}$.


\subsection{Stationarity Equations (Stage-by-Stage)}

\textbf{State stationarity} at stage $k$:
\begin{equation}\label{eq:stat_x}
  \tilde{Q}_k \Delta x_k + \tilde{N}_k \Delta u_k
  + F_k^\top \Delta \lambda_k - \Delta \lambda_{k-1} = -r_{x,k}.
\end{equation}

\textbf{Control stationarity} at stage $k$:
\begin{equation}\label{eq:stat_u}
  \tilde{N}_k^\top \Delta x_k + \tilde{R}_k \Delta u_k
  + G_k^\top \Delta \lambda_k = -r_{u,k}.
\end{equation}

\textbf{Dynamics} at stage $k$:
\begin{equation}\label{eq:dyn}
  F_k \Delta x_k + G_k \Delta u_k - \Delta x_{k+1} = -r_{p,k}.
\end{equation}


%=============================================================================
\section{Backward Riccati Recursion}
%=============================================================================

\subsection{Why Backward, Not Forward?}

In a forward Schur complement approach, we eliminate $x_0$, then $x_1$, etc.
At each stage the Schur complement is:
\begin{equation}
  S_k = H_k - C_k^\top S_{k-1}^{-1} C_k.
\end{equation}
If $S_{k-1}$ has a small eigenvalue (barely positive), then $S_{k-1}^{-1}$ has
a huge eigenvalue, and $C^\top S^{-1} C$ can be enormous. This makes $S_k$
\emph{arbitrarily negative}---one barely-PD stage causes a catastrophic
\textbf{cascade} through all subsequent stages.

The backward Riccati avoids this entirely through the \emph{floor property}
(Proposition~\ref{prop:floor}).


\subsection{The Riccati Ansatz}

We maintain the affine relation (propagating backward from terminal):
\begin{equation}\label{eq:ansatz}
  \Delta \lambda_k = P_{k+1} \Delta x_{k+1} + p_{k+1},
\end{equation}
where $P_{k+1} \in \mathbb{R}^{n_x \times n_x}$ is the symmetric cost-to-go
Hessian and $p_{k+1} \in \mathbb{R}^{n_x}$ is the cost-to-go gradient.

\textbf{Terminal condition.} At stage $N$, stationarity gives
$\Delta \lambda_{N-1} = \tilde{Q}_N \Delta x_N + r_{x,N}$, so:
\begin{equation}\label{eq:terminal}
  P_N = \tilde{Q}_N, \qquad p_N = r_{x,N}.
\end{equation}


\subsection{Deriving the Backward Recursion}

Given $P_{k+1}$ and $p_{k+1}$, we derive the update at stage $k$.

\textbf{Step 1.} Substitute the dynamics~\eqref{eq:dyn} into the ansatz~\eqref{eq:ansatz}:
\begin{equation}\label{eq:dlambda_expand}
  \Delta \lambda_k = P_{k+1} F_k \Delta x_k + P_{k+1} G_k \Delta u_k
                   + \bigl(p_{k+1} - P_{k+1} r_{p,k}\bigr).
\end{equation}

\textbf{Step 2.} Substitute into the control stationarity~\eqref{eq:stat_u}:
\begin{equation}
  \bigl[\tilde{R}_k + G_k^\top P_{k+1} G_k\bigr] \Delta u_k
  + \bigl[\tilde{N}_k^\top + G_k^\top P_{k+1} F_k\bigr] \Delta x_k = -r_{M,k},
\end{equation}
where $r_{M,k} = r_{u,k} + G_k^\top(p_{k+1} - P_{k+1} r_{p,k})$.

This defines the \textbf{control Hessian} and \textbf{cross term}:
\begin{equation}\label{eq:Mk}
  \boxed{M_k := \tilde{R}_k + G_k^\top P_{k+1} G_k} \in \mathbb{R}^{n_u \times n_u},
\end{equation}
\begin{equation}\label{eq:Lk}
  L_k := \tilde{N}_k + F_k^\top P_{k+1} G_k \in \mathbb{R}^{n_x \times n_u}.
\end{equation}

\textbf{Step 3.} Solve for the control step (requires $M_k \succ 0$):
\begin{equation}\label{eq:Du}
  \Delta u_k = -M_k^{-1}\bigl(L_k^\top \Delta x_k + r_{M,k}\bigr)
             = -K_k \Delta x_k - k_k,
\end{equation}
where the \textbf{feedback gain} and \textbf{feedforward} are:
\begin{equation}
  K_k = M_k^{-1} L_k^\top \in \mathbb{R}^{n_u \times n_x}, \qquad
  k_k = M_k^{-1} r_{M,k} \in \mathbb{R}^{n_u}.
\end{equation}

\textbf{Step 4.} Substitute back into state stationarity to get the Riccati update:
\begin{equation}\label{eq:riccati}
  \boxed{P_k = \tilde{Q}_k + F_k^\top P_{k+1} F_k - L_k M_k^{-1} L_k^\top.}
\end{equation}


\subsection{The Anti-Cascade Floor Property}

\begin{proposition}[Floor Property]\label{prop:floor}
If $M_k \succ 0$ and $\tilde{N}_k = 0$, then
\begin{equation}
  P_k \succeq \tilde{Q}_k.
\end{equation}
\end{proposition}

\begin{proof}
With $\tilde{N}_k = 0$, we have $L_k = F_k^\top P_{k+1} G_k$ and
\begin{equation}
  P_k = \tilde{Q}_k + F_k^\top\!\bigl[P_{k+1}
        - P_{k+1} G_k (G_k^\top P_{k+1} G_k + \tilde{R}_k)^{-1} G_k^\top P_{k+1}\bigr] F_k.
\end{equation}
The bracketed matrix is the Schur complement of $M_k$ in
\begin{equation}
  \Omega = \begin{bmatrix}
    P_{k+1} & P_{k+1} G_k \\
    G_k^\top P_{k+1} & G_k^\top P_{k+1} G_k + \tilde{R}_k
  \end{bmatrix}
  = \begin{bmatrix} I \\ G_k^\top \end{bmatrix} P_{k+1}
    \begin{bmatrix} I & G_k \end{bmatrix}
  + \blockdiag(0, \tilde{R}_k).
\end{equation}
Since $P_{k+1} \succeq 0$ (by induction) and $\tilde{R}_k \succeq 0$ (after BISC),
$\Omega \succeq 0$. The Schur complement of a PSD block in a PSD matrix is PSD.
Therefore $P_k = \tilde{Q}_k + F_k^\top(\text{PSD})F_k \succeq \tilde{Q}_k$.
\end{proof}

\textbf{Consequence:} The backward Riccati \emph{cannot amplify indefiniteness}.
If $\tilde{Q}_k \succ 0$ (which the barrier contribution helps ensure), then
$P_k$ is automatically PD. No cascade.

\begin{remark}
For general $\tilde{N}_k \ne 0$:
$P_k \succeq \tilde{Q}_k - \tilde{N}_k M_k^{-1} \tilde{N}_k^\top$,
which is the Schur complement of $M_k$ in $\tilde{H}_k$.
\end{remark}


\subsection{Forward Sweep (Step Recovery)}

After the backward pass computes $P_k$, $K_k$, $k_k$ for all $k$, the Newton
step is recovered forward:
\begin{equation}
\begin{aligned}
  &\Delta x_0 \text{ from initial condition}, \\
  &\text{For } k = 0, \dots, N\!-\!1\text{:} \\
  &\quad \Delta u_k = -K_k \Delta x_k - k_k, \\
  &\quad \Delta x_{k+1} = F_k \Delta x_k + G_k \Delta u_k - r_{p,k}, \\
  &\quad \Delta \lambda_k = P_{k+1} \Delta x_{k+1} + p_{k+1}.
\end{aligned}
\end{equation}


%=============================================================================
\section{The BISC Algorithm}
%=============================================================================

\subsection{Where Indefiniteness Enters}

The backward Riccati requires inverting $M_k$ at each stage. From~\eqref{eq:Mk}:
\begin{equation}
  M_k = \underbrace{\tilde{R}_k}_{\text{can be indefinite}} + \underbrace{G_k^\top P_{k+1} G_k}_{\succeq 0}.
\end{equation}
The term $\tilde{R}_k$ contains the dynamics curvature $\lambda_k^\top \nabla^2_{u_k u_k} f_k$
from AD, which can be negative. When this negative curvature is large enough,
$M_k$ becomes indefinite and the Riccati fails. \textbf{This is where BISC intervenes.}


\subsection{PD Cone Projection (Minimum Frobenius-Norm)}

\begin{definition}[PD Cone Projection]
For a symmetric $A \in \mathbb{R}^{m \times m}$ with eigendecomposition
$A = Q \Lambda Q^\top$ and threshold $\delta > 0$:
\begin{equation}\label{eq:proj}
  \proj_\delta(A) := Q\,\max(\Lambda, \delta I)\,Q^\top,
\end{equation}
where $\max$ is applied element-wise to the eigenvalue diagonal.
\end{definition}

\begin{lemma}[Optimality of PD Projection]\label{lem:proj}
The modification $E = \proj_\delta(A) - A$ satisfies
\begin{equation}
  E = Q\,\max(\delta I - \Lambda, 0)\,Q^\top \succeq 0,
\end{equation}
and is the \emph{unique minimizer} of
\begin{equation}
  \min_{E \succeq 0} \|E\|_F \quad \text{subject to} \quad
  \lambda_{\min}(A + E) \ge \delta.
\end{equation}
\end{lemma}

\begin{proof}
In the eigenbasis of $A$, the Frobenius norm decomposes:
$\|E\|_F^2 = \sum_i e_i^2$, with constraints $\lambda_i + e_i \ge \delta$,
$e_i \ge 0$. The minimum is $e_i = \max(\delta - \lambda_i, 0)$ independently
for each $i$.
\end{proof}

\textbf{Key properties:}
\begin{enumerate}
  \item Only eigenvalues $< \delta$ are modified; healthy eigenvalues are untouched.
  \item If $A \succeq \delta I$, then $E = 0$ (no modification).
  \item For $n_u \times n_u$ blocks (typically $n_u = 1$ to $10$), the
        eigendecomposition costs $O(n_u^3)$---negligible.
\end{enumerate}


\subsection{The Cheng--Higham Modified Cholesky}

For implementation, especially with larger control dimensions, we use the
Cheng--Higham (1998) modified Cholesky based on symmetric indefinite factorization.

\textbf{Background.} Given a symmetric $A \in \mathbb{R}^{m \times m}$
(not necessarily PD), the Bunch--Kaufman algorithm computes:
\begin{equation}\label{eq:lbl}
  PAP^\top = LBL^\top,
\end{equation}
where $P$ is a permutation (pivoting), $L$ is unit lower triangular, and $B$ is
block-diagonal with $1 \times 1$ and $2 \times 2$ blocks.

\begin{algorithm}[H]
\caption{Cheng--Higham Modified Cholesky}
\label{alg:chenghigham}
\begin{algorithmic}[1]
\REQUIRE Symmetric $A \in \mathbb{R}^{m \times m}$, threshold $\delta > 0$
\ENSURE Modified $\tilde{A} = A + E$ with $\tilde{A} \succeq \delta I$, $\|E\|_F$ minimized
\STATE Compute symmetric indefinite factorization $PAP^\top = LBL^\top$
       \COMMENT{Bunch--Kaufman}
\FOR{each diagonal block $B_j$ of $B$}
  \IF{$B_j$ is $1 \times 1$, i.e., $B_j = [b_{jj}]$}
    \IF{$b_{jj} < \delta$}
      \STATE $\tilde{b}_{jj} \leftarrow \delta$
             \COMMENT{Clip to threshold}
    \ENDIF
  \ELSIF{$B_j$ is $2 \times 2$}
    \STATE Compute eigenvalues of $B_j$:
    \STATE \quad $\tau \leftarrow (b_{jj} + b_{j+1,j+1})/2$
    \STATE \quad $\rho \leftarrow \sqrt{((b_{jj} - b_{j+1,j+1})/2)^2 + b_{j,j+1}^2}$
    \STATE \quad $\lambda_1 \leftarrow \tau + \rho$, \quad
           $\lambda_2 \leftarrow \tau - \rho$
    \IF{$\min(\lambda_1, \lambda_2) < \delta$}
      \STATE $\tilde{\lambda}_i \leftarrow \max(\lambda_i, \delta)$ for $i = 1, 2$
      \STATE Compute eigenvectors $V$ of $B_j$
      \STATE $\tilde{B}_j \leftarrow V \diag(\tilde{\lambda}_1, \tilde{\lambda}_2) V^\top$
    \ENDIF
  \ENDIF
\ENDFOR
\STATE Reconstruct: $\tilde{A} \leftarrow P^\top L \tilde{B} L^\top P$
\end{algorithmic}
\end{algorithm}

\textbf{Equivalence for small blocks.} For the $n_u \times n_u$ control Hessian
$M_k$ (typically $n_u = 1$ to $10$), the Bunch--Kaufman factorization captures
all eigenvalues through its $1 \times 1$ and $2 \times 2$ diagonal blocks.
Clipping the eigenvalues of each $B_j$ is equivalent to
$\proj_\delta(M_k)$---the same minimum Frobenius-norm projection. For very
small blocks ($n_u \le 4$), direct eigendecomposition is simpler; for larger
blocks ($n_u > 4$), Cheng--Higham is more numerically stable.


\subsection{Choice of $\delta$ (No Free Parameters)}

\begin{equation}\label{eq:delta}
  \delta = \kappa \cdot \mu, \qquad \kappa > 0 \text{ fixed (e.g., } \kappa = 0.1\text{)}.
\end{equation}

\textbf{Rationale:}
\begin{itemize}
  \item Early iterations ($\mu$ large): $\delta$ is large, providing robust regularization.
  \item Near solution ($\mu$ small): $\delta$ is small, preserving Newton step quality.
  \item At convergence ($\mu \to 0$): $\delta \to 0$, and under SOSC the modification
        vanishes completely (Theorem~\ref{thm:sosc}).
\end{itemize}
No trial-and-error. No multiple factorizations. The parameter is determined by
the IPM's own barrier parameter.


\subsection{The Complete BISC Algorithm}

\begin{algorithm}[H]
\caption{BISC: Block Inertia-free Structural Convexification}
\label{alg:bisc}
\begin{algorithmic}[1]
\REQUIRE Stage Hessians $\{\tilde{Q}_k, \tilde{R}_k, \tilde{N}_k\}$, Jacobians $\{F_k, G_k\}$,
         residuals $\{r_{x,k}, r_{u,k}, r_{p,k}\}$, barrier parameter $\mu$
\STATE $\delta \leftarrow \kappa \cdot \mu$
\STATE \textbf{--- Backward Pass ---}
\STATE $\tilde{P}_N \leftarrow \proj_\delta(\tilde{Q}_N)$; \quad $p_N \leftarrow r_{x,N}$
\FOR{$k = N\!-\!1, \dots, 0$}
  \STATE $M_k \leftarrow \tilde{R}_k + G_k^\top \tilde{P}_{k+1} G_k$
         \COMMENT{Control Hessian ($n_u \times n_u$)}
  \STATE $\tilde{M}_k \leftarrow \proj_\delta(M_k)$
         \COMMENT{\textbf{BISC core: min-norm PD projection}}
  \STATE $L_k \leftarrow \tilde{N}_k + F_k^\top \tilde{P}_{k+1} G_k$
  \STATE $K_k \leftarrow \tilde{M}_k^{-1} L_k^\top$; \quad
         $k_k \leftarrow \tilde{M}_k^{-1} r_{M,k}$
         \COMMENT{Solve via Cholesky of $\tilde{M}_k$}
  \STATE $P_k \leftarrow \tilde{Q}_k + F_k^\top \tilde{P}_{k+1} F_k - L_k K_k$
         \COMMENT{Riccati update}
  \STATE $\tilde{P}_k \leftarrow \proj_\delta(P_k)$
         \COMMENT{Safety projection (often $\tilde{P}_k = P_k$)}
  \STATE Compute $p_k$ from residuals
\ENDFOR
\STATE \textbf{--- Forward Sweep ---}
\STATE Compute $\Delta x_0$ from initial condition
\FOR{$k = 0, \dots, N\!-\!1$}
  \STATE $\Delta u_k \leftarrow -K_k \Delta x_k - k_k$
  \STATE $\Delta x_{k+1} \leftarrow F_k \Delta x_k + G_k \Delta u_k - r_{p,k}$
  \STATE $\Delta \lambda_k \leftarrow \tilde{P}_{k+1} \Delta x_{k+1} + p_{k+1}$
\ENDFOR
\STATE Recover slack and inequality multiplier steps from eliminated equations
\end{algorithmic}
\end{algorithm}

\textbf{Computational cost:}
$O(N(n_x + n_u)^3)$ for the backward pass, $O(N(n_x + n_u)^2)$ for the forward
sweep. \emph{Linear} in horizon $N$. One backward + one forward pass. No retries.


%=============================================================================
\section{Theoretical Analysis}
%=============================================================================

\subsection{Theorem 1: Minimum Perturbation}

\begin{theorem}[Stagewise Minimum Frobenius-Norm Perturbation]\label{thm:minpert}
At each stage $k$, the BISC modification $E_k^M = \tilde{M}_k - M_k$ satisfies
\begin{equation}
  \|E_k^M\|_F = \min\bigl\{\|E\|_F : E \succeq 0, \;
  \lambda_{\min}(M_k + E) \ge \delta\bigr\}.
\end{equation}
\end{theorem}

\begin{proof}
Direct application of Lemma~\ref{lem:proj} to $M_k$ with threshold $\delta$.
\end{proof}

\begin{corollary}[Structured Sparsity]
The total BISC modification satisfies
\begin{equation}
  \|E_{\mathrm{total}}\|_F^2 = \sum_k \|E_k^M\|_F^2 + \sum_k \|E_k^P\|_F^2.
\end{equation}
Each $E_k^M \in \mathbb{R}^{n_u \times n_u}$ and $E_k^P \in \mathbb{R}^{n_x \times n_x}$.
Off-diagonal blocks and the constraint Jacobian $J$ are \emph{never modified}.
\end{corollary}


\subsection{Theorem 2: Descent Direction}

\begin{theorem}[Descent Guarantee]\label{thm:descent}
If $\tilde{P}_k \succ 0$ for all $k = 0, \dots, N$ (ensured by Algorithm~\ref{alg:bisc}),
then the computed step $(\Delta z, \Delta \lambda)$ satisfies
\begin{equation}
  \Delta z^\top \nabla \varphi_\mu(z) < 0,
\end{equation}
where $\varphi_\mu(z) = \sum_k \ell_k(z_k) - \mu \sum_k \mathbf{1}^\top \log(-g_k(z_k))$
is the log-barrier merit function.
\end{theorem}

\begin{proof}
The BISC step solves the modified KKT system
\begin{equation}
  \begin{bmatrix} \tilde{H} + E & J^\top \\ J & 0 \end{bmatrix}
  \begin{bmatrix} \Delta z \\ \Delta \lambda \end{bmatrix}
  = -\begin{bmatrix} \nabla \varphi_\mu \\ c \end{bmatrix},
\end{equation}
where $E = \blockdiag(E_0, \dots, E_N)$ is the total BISC modification.

Let $Z$ span $\ker(J)$. The reduced Hessian is
$Z^\top(\tilde{H} + E)Z$. The backward Riccati with BISC computes an implicit
$LDL^\top$ factorization of this reduced Hessian, with all diagonal blocks
($\tilde{M}_k$ and $\tilde{P}_k$) PD by construction. Therefore
$Z^\top(\tilde{H} + E)Z \succ 0$.

By standard IPM theory (Nocedal \& Wright, Ch.~19), when the reduced Hessian
is PD:
$\Delta z^\top \nabla \varphi_\mu = -\Delta z^\top(\tilde{H} + E)\Delta z < 0$.
\end{proof}


\subsection{Theorem 3: BISC vs.\ IPOPT}

\begin{theorem}[Smaller Perturbation than IPOPT]\label{thm:ipopt}
Let $\delta_{\mathrm{IP}}$ be IPOPT's inertia correction
($\tilde{H} \leftarrow \tilde{H} + \delta_{\mathrm{IP}} I$). Then
\begin{equation}
  \|E_{\mathrm{BISC}}\|_F \le \|\delta_{\mathrm{IP}} I_{n_{\mathrm{total}}}\|_F
  = \delta_{\mathrm{IP}} \sqrt{n_{\mathrm{total}}},
\end{equation}
with strict inequality whenever any stage has $E_k^M = 0$ or $E_k^P = 0$.
\end{theorem}

\begin{proof}
IPOPT adds $\delta_{\mathrm{IP}}$ to all $n_{\mathrm{total}} = \sum_k n_{z,k}$
diagonal elements. BISC adds $E_k^M$ only where $M_k$ is indefinite and $E_k^P$
only where $P_k$ is indefinite:
\begin{equation}
  \|E_{\mathrm{BISC}}\|_F^2 = \sum_{k \in \mathcal{S}_{\mathrm{bad}}} \|E_k\|_F^2,
\end{equation}
where $\mathcal{S}_{\mathrm{bad}} \subseteq \{0, \dots, N\}$. Since typically
$|\mathcal{S}_{\mathrm{bad}}| \ll N+1$, the inequality is strict.
\end{proof}


\subsection{Theorem 4: SOSC Recovery and Quadratic Convergence}

\begin{definition}[SOSC]
The second-order sufficient condition holds at $(z^*, \lambda^*, \nu^*)$ if
\begin{equation}
  d^\top \nabla^2_{zz} \mathcal{L}\, d > 0 \quad
  \text{for all } d \ne 0 \text{ with } Jd = 0 \text{ and active constraint conditions}.
\end{equation}
Let $\sigma > 0$ denote the minimum reduced Hessian eigenvalue at the solution.
\end{definition}

\begin{theorem}[Exact Newton Recovery]\label{thm:sosc}
Suppose SOSC holds with minimum eigenvalue $\sigma > 0$. There exists a
neighborhood of $(z^*, \lambda^*)$ such that for all iterates in this neighborhood:
\begin{enumerate}
  \item[(i)] $\lambda_{\min}(M_k) \ge \sigma / c$ for some constant $c > 0$, at all stages.
  \item[(ii)] For $\delta < \sigma / c$ (which holds for $\mu$ sufficiently small):
              $E_k^M = 0$ at all stages.
  \item[(iii)] $E_k^P = 0$ at all stages.
  \item[(iv)] BISC produces the \emph{exact} Newton step (zero modification).
\end{enumerate}
\end{theorem}

\begin{proof}
\emph{(i)} Under SOSC, the reduced Hessian eigenvalues are continuous and bounded
below by $\sigma$ at the solution. By continuity, they are bounded below by
$\sigma/2 > 0$ in a neighborhood. The control Hessian $M_k$ is related to the
reduced Hessian through the Riccati factorization; its eigenvalues are bounded
below by $\sigma/c$ for some $c > 0$ (eigenvalue interlacing).

\emph{(ii)} With $\delta = \kappa \mu \to 0$, eventually $\delta < \sigma/c$.
Then $\lambda_{\min}(M_k) \ge \sigma/c > \delta$, so $\proj_\delta(M_k) = M_k$.

\emph{(iii)} By Proposition~\ref{prop:floor}, $P_k \succeq \tilde{Q}_k$.
Barrier contributions ensure $\tilde{Q}_k \succ 0$ near the solution.

\emph{(iv)} All modifications are zero; BISC computes the exact Newton step.
\end{proof}

\begin{corollary}[Quadratic Convergence with AD]\label{cor:quadratic}
Since BISC eventually gives exact Newton steps and AD provides exact Hessians:
\begin{equation}
  \|(z_{j+1}, \lambda_{j+1}) - (z^*, \lambda^*)\|
  = O\!\left(\|(z_j, \lambda_j) - (z^*, \lambda^*)\|^2\right).
\end{equation}
Other methods (IPOPT's $\delta_{\mathrm{IP}} > 0$, Vanroye's $\rho_d > 0$)
always have nonzero perturbation and therefore cap convergence at superlinear.
\end{corollary}


\subsection{Comparison with Existing Methods}

\begin{table}[h]
\centering
\caption{BISC vs.\ Verschueren et al.\ (2017)}
\begin{tabular}{lll}
\toprule
\textbf{Aspect} & \textbf{Verschueren} & \textbf{BISC} \\
\midrule
Context & SQP (QP subproblems) & IPM (barrier subproblem) \\
Projection & PROJECT / MIRROR heuristic & Min Frobenius-norm \\
Parameter & $\varepsilon$ (free, user-chosen) & $\delta = \kappa\mu$ (from barrier) \\
Barrier benefit & N/A & Reduces modification needed \\
Theory & SQP convergence & SOSC exact recovery \\
\bottomrule
\end{tabular}
\end{table}

\begin{table}[h]
\centering
\caption{BISC vs.\ Vanroye et al.\ (2025)}
\begin{tabular}{lll}
\toprule
\textbf{Aspect} & \textbf{Vanroye} & \textbf{BISC} \\
\midrule
Modification & Dual ($-\rho_d I$ in constraint block) & Primal ($M_k$ block) \\
Selectivity & Global (all stages) & Stage-selective \\
Constraints & Modified ($J^\top$ structure changed) & Preserved \\
Norm optimality & Scalar $\rho_d$ (over-regularizes) & Frobenius-optimal per block \\
Near solution & Always $\rho_d > 0$ & $E = 0$ under SOSC \\
Convergence & Superlinear (at best) & Quadratic (AD + exact Newton) \\
\bottomrule
\end{tabular}
\end{table}


%=============================================================================
\section{The Role of Automatic Differentiation}
%=============================================================================

\subsection{Why AD + BISC}

The Lagrangian Hessian~\eqref{eq:Wk} has three terms:
\begin{enumerate}
  \item[\textbf{I.}] $\nabla^2 \ell_k$ --- objective curvature (usually PSD).
  \item[\textbf{II.}] $\lambda_k^\top \nabla^2 f_k$ --- dynamics curvature (\emph{can be negative}).
  \item[\textbf{III.}] $\nu_k^\top \nabla^2 g_k$ --- inequality curvature.
\end{enumerate}

\textbf{Gauss--Newton} uses only term~I: PSD, no convexification needed, but
\emph{linear} convergence. \textbf{AD exact Hessians} include all three terms:
\emph{quadratic} convergence, but indefiniteness. BISC handles the indefiniteness
with minimum perturbation.

\begin{table}[h]
\centering
\caption{Convergence hierarchy}
\begin{tabular}{llll}
\toprule
\textbf{Method} & \textbf{Hessian} & \textbf{Convergence} & \textbf{Convexification} \\
\midrule
Gauss--Newton         & Term I only   & Linear       & Not needed \\
L-BFGS                & Approximate   & Superlinear  & Not needed (PD by construction) \\
Exact + IPOPT         & All (AD)      & Superlinear  & Global $\delta I$ (wasteful) \\
Exact + BISC          & All (AD)      & Quadratic    & Min-norm per block \\
\bottomrule
\end{tabular}
\end{table}


\subsection{What AD Provides to BISC}

At each stage, AD computes:
\begin{itemize}
  \item First derivatives: $F_k$, $G_k$, $J_{g,k}$ (Jacobians).
  \item Second derivatives: $\tilde{Q}_k$, $\tilde{R}_k$, $\tilde{N}_k$ (Hessian blocks).
\end{itemize}
The AD computation graph automatically reveals the block-diagonal structure
of $\tilde{H}$ (since $f_k$ depends only on $z_k$), enabling automatic
detection of OC structure for the Riccati-based BISC algorithm.


%=============================================================================
\section{Summary of Novelty}
%=============================================================================

BISC is novel on five axes compared to all existing methods:
\begin{enumerate}
  \item \textbf{IPM context} --- unlike Verschueren et al.\ (SQP only).
  \item \textbf{Primal modification} --- unlike Vanroye et al.\ (dual regularization).
  \item \textbf{Stage-selective} --- unlike IPOPT (global) and Vanroye (all stages).
  \item \textbf{Min Frobenius-norm} --- unlike IPOPT (scalar) and Vanroye (scalar $\rho_d$).
  \item \textbf{SOSC exact recovery} --- unlike IPOPT ($\delta > 0$ always)
        and Vanroye ($\rho_d > 0$ always).
\end{enumerate}

\textbf{One sentence:} BISC applies the minimum-norm PSD projection to only the
small control Hessian blocks that need it during backward Riccati, tied to the
barrier parameter, giving guaranteed descent with zero modification near the solution.


%=============================================================================
\section{Numerical Experiments (Planned)}
%=============================================================================

\begin{enumerate}
  \item \textbf{Space shuttle reentry} ($n_x=6$, $n_u=2$, $N=100$): Highly nonlinear
        dynamics. Compare BISC vs IPOPT vs Gauss--Newton.
  \item \textbf{Race car lap optimization} ($n_x=6\text{--}8$, $n_u=2\text{--}4$, $N=200\text{--}500$):
        Path constraints. Show stage-selectivity.
  \item \textbf{Convergence rates}: iteration counts, trial factorization counts.
  \item \textbf{Perturbation norms}: $\|E_{\mathrm{BISC}}\|_F$ vs
        $\|\delta_{\mathrm{IP}} I\|_F$ across IPM iterations.
  \item \textbf{Timing}: wall-clock comparison showing single-pass vs retries.
\end{enumerate}


\bibliographystyle{siam}
\begin{thebibliography}{10}

\bibitem{cheng1998}
S.~H. Cheng and N.~J. Higham,
\emph{A modified Cholesky algorithm based on a symmetric indefinite factorization},
SIAM J. Matrix Anal. Appl., 19 (1998), pp.~1097--1110.

\bibitem{verschueren2017}
R.~Verschueren, M.~Zanon, R.~Quirynen, and M.~Diehl,
\emph{A sparsity preserving convexification procedure for indefinite QP arising in direct optimal control},
Numerical Algebra, Control and Optimization, 7 (2017), pp.~365--379.

\bibitem{vanroye2025}
M.~Vanroye, et al.,
\emph{Dual-regularized Riccati recursions for interior point methods},
arXiv:2509.16370, 2025.

\bibitem{rao1998}
C.~V. Rao, S.~J. Wright, and J.~B. Rawlings,
\emph{Application of interior-point methods to model predictive control},
J. Optim. Theory Appl., 99 (1998), pp.~723--757.

\bibitem{nocedal2006}
J.~Nocedal and S.~J. Wright,
\emph{Numerical Optimization}, 2nd ed., Springer, 2006.

\bibitem{frison2020}
G.~Frison and M.~Diehl,
\emph{HPIPM: a high-performance quadratic programming framework for model predictive control},
IFAC-PapersOnLine, 53 (2020), pp.~6563--6569.

\bibitem{wachter2006}
A.~W\"achter and L.~T. Biegler,
\emph{On the implementation of an interior-point filter line-search algorithm for large-scale nonlinear programming},
Math. Programming, 106 (2006), pp.~25--57.

\bibitem{cartis2011}
C.~Cartis, N.~I.~M. Gould, and Ph.~L. Toint,
\emph{Adaptive cubic regularisation methods for unconstrained optimization. Part I: motivation, convergence and numerical results},
Math. Programming, 127 (2011), pp.~245--295.

\bibitem{dauphin2014}
Y.~Dauphin, R.~Pascanu, C.~Gulcehre, K.~Cho, S.~Ganguli, and Y.~Bengio,
\emph{Identifying and attacking the saddle point problem in high-dimensional non-convex optimization},
NeurIPS, 2014.

\end{thebibliography}

\end{document}
